\documentclass{article}
\usepackage{graphicx}
\usepackage{amsmath}
\usepackage{setspace}
\onehalfspacing
\usepackage{tikz}
\usepackage{pgfplots}
\pgfplotsset{compat=1.18}

\title{CS556 - Midterm Retake}
\author{John Rizzo}
\date{November 19, 2025}

\usepackage[margin=2cm]{geometry}

\begin{document}

\maketitle

%
% Question 1 %%%%%%%%%%%%%%%%%%%%%%%%%%%%%%%%%%%%%%%%%%%%%%%%%%%%%%%%%%%%%%%%%%
\section*{Question 1} (20 points): Your answer to any of the following questions should be 'in-sufficient information' if you feel the question does not supply enough information to determine the correct answer.  Your answer can also be 'does not exist' if you judge that the question is about a property that is not applicable to the matrix in question.  Given that a square matrix of dimension 45 has a linearly independent set of vectors (i.e., it is full rank), answer the following questions.
\begin{enumerate}
    \item What is the rank of the matrix? \\
          The rank of the matrix is 45, since it is given that the matrix is full rank.
    \item What is the numerical value of the determinant of the matrix? \\
          The determinant of a full rank square matrix is non-zero. However, without specific numerical values, we cannot determine the exact value of the determinant. Therefore, the answer is 'in-sufficient information'.
    \item What is the rank of the column space? \\
          The rank of the column space is equal to the rank of the matrix, which is 45.
    \item What is the dimension of the null space (i.e. the nullity)? \\
            The nullity of the matrix can be calculated using the formula: 
            $ \text{nullity} = \text{number of columns} - \text{rank}$
            Since the matrix is square with 45 columns and has a rank of 45, the nullity is:
            $ \text{nullity} = 45 - 45 = 0$
    \item What is the dimension of the left null space (i.e. the nullity of the left null space)? \\
            The dimension of the left null space can be calculated using the formula:
            $\text{left nullity} = \text{number of rows} - \text{rank}$
            Since the matrix is square with 45 rows and has a rank of 45, the left nullity is:
            $\text{left nullity} = 45 - 45 = 0$
    \item Is the matrix invertible (why or why not)? \\
            The matrix is invertible because it is square and full rank, which means its determinant is non-zero.
\end{enumerate}

%
% Question 2 %%%%%%%%%%%%%%%%%%%%%%%%%%%%%%%%%%%%%%%%%%%%%%%%%%%%%%%%%%%%%%%%%
\vfill\break
\section*{Question 2} (20 points): Calculate the solution (if it exists) to the following system of linear equations (SLE) and demonstrate your methodology.

\noindent
$x_1 - \frac{4}{7}x_2 + \frac{2}{7}x_3 = \frac{11}{7}$ \\
$-\frac{2}{7}x_2 + \frac{1}{7}x_3 = \frac{2}{7}$ \\
$\frac{9}{7}x_2 + \frac{6}{7}x_3 = \frac{33}{7}$

\noindent
\textbf{Solution}:

$\begin{bmatrix}
1 & -\frac{4}{7} & \frac{2}{7} & | & \frac{11}{7} \\
0 & -\frac{2}{7} & \frac{1}{7} & | & \frac{2}{7} \\
0 & \frac{9}{7} & \frac{6}{7} & | & \frac{33}{7}
\end{bmatrix}$

\noindent
$R2 = R2 + R3$
$\begin{bmatrix}
1 & -\frac{4}{7} & \frac{2}{7} & | & \frac{11}{7} \\
0 & 1            & 1           & | & 5 \\
0 & \frac{9}{7}  & \frac{6}{7} & | & \frac{33}{7}
\end{bmatrix}$

\noindent
$R3 = R3 \times 7$
$\begin{bmatrix}
1 & -\frac{4}{7} & \frac{2}{7} & | & \frac{11}{7} \\
0 & 1            & 1           & | & 5 \\
0 & 9            & 6           & | & 33
\end{bmatrix}$

\noindent
$R3 = R3 - 9R2$
$\begin{bmatrix}
1 & -\frac{4}{7} & \frac{2}{7} & | & \frac{11}{7} \\
0 & 1            & 1           & | & 5 \\
0 & 0            & -3          & | & -12
\end{bmatrix}$

\noindent
$R3 = R3 \times -\frac{1}{3}$
$\begin{bmatrix}
1 & -\frac{4}{7} & \frac{2}{7} & | & \frac{11}{7} \\
0 & 1            & 1           & | & 5 \\
0 & 0            & 1           & | & 4
\end{bmatrix}$

\noindent
$R1 = \frac{4}{7}R2 + R1$
$\begin{bmatrix}
1 & 0 & \frac{6}{7} & | & \frac{31}{7} \\
0 & 1 & 1           & | & 5 \\
0 & 0 & 1           & | & 4
\end{bmatrix}$

\noindent
$R2 = R2 - R3$
$\begin{bmatrix}
1 & 0 & \frac{6}{7} & | & \frac{31}{7} \\
0 & 1 & 0           & | & 1 \\
0 & 0 & 1           & | & 4
\end{bmatrix}$

\noindent
$R1 = R1 - \frac{6}{7}R3$
$\begin{bmatrix}
1 & 0 & 0 & | & 1 \\
0 & 1 & 0 & | & 1 \\
0 & 0 & 1 & | & 4
\end{bmatrix}$

%
% Question 3 %%%%%%%%%%%%%%%%%%%%%%%%%%%%%%%%%%%%%%%%%%%%%%%%%%%%%%%%%%%%%%%%%
\vfill\break
\section*{Question 3} (15 points): Answer whether the following statements are True, False, or if there is insufficient information to make the determination.  If $a$, $b$, and $c$ are a set of linearly independent vectors, each in $R^3$ and none of which are the zero vector, which of the following statements are true about them?

\begin{enumerate}
    \item $a$, $b$, and $c$ span the space $R^3$. \\
          Since $a$, $b$, and $c$ are linearly independent vectors in $R^3$ and there are three of them, they do span the space $R^3$. \textbf{True}
    \item There exists a set of scalars $\lambda_1$, $\lambda_2$, and $\lambda_3$ where $\lambda_1 \neq \lambda_2$ which can be combined to form $\lambda_1 a + \lambda_2 b + \lambda_3 c = 0$. \\
          \textbf{False}. Since $a$, $b$, and $c$ are linearly independent, the only solution to the equation $\lambda_1 a + \lambda_2 b + \lambda_3 c = 0$ is when all scalars are zero.
    \item The determinant of the matrix with $a$, $b$, and $c$ as columns is 0. \\
          \textbf{False}. The determinant of a matrix with linearly independent columns is non-zero.
    \item Consider a specific instantiation of the vectors as $a = \begin{bmatrix}
        1 \\ 2 \\ 3
    \end{bmatrix}$, $b = \begin{bmatrix}
        -1 \\ 4 \\ 9
    \end{bmatrix}$, and $c = \begin{bmatrix}
        -11 \\ 10 \\ 21
    \end{bmatrix}$. Could this represent the set of vectors described in the original question?
        To determine if these vectors are linearly independent, we can form a matrix with these vectors as columns and calculate its determinant:
        \[M = \begin{bmatrix}
        1 & -1 & -11 \\
        2 & 4 & 10 \\
        3 & 9 & 21
        \end{bmatrix}\]
        The determinant of matrix $M$ is calculated as follows:
        \[
        \text{det}(M) = 1 \cdot (4 \cdot 21 - 9 \cdot 10) - (-1) \cdot (2 \cdot 21 - 10 \cdot 3) + (-11) \cdot (2 \cdot 9 - 4 \cdot 3)
        \]
        \[
        \text{det}(M) = 1 \cdot (84 - 90) - (-1) \cdot (42 - 30) + (-11) \cdot (18 - 12)
        \]
        \[
        \text{det}(M) = 1 \cdot (-6) - (-1) \cdot 12 + (-11) \cdot 6
        \]
        \[
        \text{det}(M) = -6 + 12 - 66 = -60
        \]
        Since the determinant is non-zero ($-60 \neq 0$), the vectors $a$, $b$, and $c$ are linearly independent. Therefore, this set of vectors could represent the set described in the original question. \textbf{True}
\end{enumerate}

%
% Question 4 %%%%%%%%%%%%%%%%%%%%%%%%%%%%%%%%%%%%%%%%%%%%%%%%%%%%%%%%%%%%%%%%%
\vfill\break
\section*{Question 4} (10 points): Calculate the inverse of 2x2 matrix $A = \begin{bmatrix}-6 & 2 \\ 9 & -3\end{bmatrix}$.

\noindent
$A^{-1} = \frac{1}{det(A)}\begin{bmatrix}a_{22} & -a_{12} \\ -a_{21} & a_{11} \end{bmatrix}$ \\
$det(A) = -6 \times -3 - 2 \times 9 = 0$ \\
The system of equations represented by the augmented matrix has no solution. Therefore, the inverse of matrix A does not exist.

%
% Question 5 %%%%%%%%%%%%%%%%%%%%%%%%%%%%%%%%%%%%%%%%%%%%%%%%%%%%%%%%%%%%%%%%%
\vfill\break
\section*{Question 5} (15 points): Employ Gram-schmidt orthogonalization to derive an orthonormal set of vectors from the set of independent vectors: $a = \begin{bmatrix}1 \\ -1 \\ 1\end{bmatrix}$, $b = \begin{bmatrix}1 \\ 0 \\ 1\end{bmatrix}$, $c = \begin{bmatrix}1 \\ 1 \\ 2\end{bmatrix}$.

\noindent
Let $A = a$\\
$A^T = \begin{bmatrix}1 & -1 & 1\end{bmatrix}$\\
$AA^T = \begin{bmatrix}1 & -1 & 1\end{bmatrix}\begin{bmatrix}1 \\ -1 \\ 1\end{bmatrix} = 3$ \\
$A' = \frac{A}{||A||} = \frac{\begin{bmatrix}1 \\ -1 \\ 1\end{bmatrix}}{\sqrt{1^2 + (-1)^2 + 1^2}} = \frac{\begin{bmatrix}1 \\ -1 \\ 1\end{bmatrix}}{\sqrt{3}} = \frac{1}{\sqrt{3}}\begin{bmatrix}1 \\ -1 \\ 1\end{bmatrix}$

\noindent
$B = b - \frac{A^Tb}{A^TA}A = \begin{bmatrix}1 \\ 0 \\ 1\end{bmatrix} - \frac{2}{3}\begin{bmatrix}1 \\ -1 \\ 1\end{bmatrix} = \begin{bmatrix}1/3 \\ 2/3 \\ 1/3\end{bmatrix}$ \\
$B^T = \begin{bmatrix}1/3 & 2/3 & 1/3\end{bmatrix}$ \\
$BB^T = \begin{bmatrix}1/3 & 2/3 & 1/3\end{bmatrix}\begin{bmatrix}1/3 \\ 2/3 \\ 1/3\end{bmatrix} = 2/3$ \\
$B' = \frac{B}{||B||} = \frac{\begin{bmatrix}1/3 \\ 2/3 \\ 1/3\end{bmatrix}}{\sqrt{2/3}} = \sqrt{\frac{3}{2}}\begin{bmatrix}1/3 \\ 2/3 \\ 1/3\end{bmatrix} = \begin{bmatrix}1/\sqrt{6} \\ 2/\sqrt{6} \\ 1/\sqrt{6}\end{bmatrix}$

\noindent
$C = c - \frac{A^Tc}{A^TA}A - \frac{B^Tc}{B^TB}B$ \\
$A^Tc = \begin{bmatrix}1 & -1 & 1\end{bmatrix}\begin{bmatrix}1 \\ 1 \\ 2\end{bmatrix} = 2$, \quad $B^Tc = \begin{bmatrix}1/3 & 2/3 & 1/3\end{bmatrix}\begin{bmatrix}1 \\ 1 \\ 2\end{bmatrix} = \frac{5}{3}$ \\
$C = \begin{bmatrix}1 \\ 1 \\ 2\end{bmatrix} - \frac{2}{3}\begin{bmatrix}1 \\ -1 \\ 1\end{bmatrix} - \frac{5/3}{2/3}\begin{bmatrix}1/3 \\ 2/3 \\ 1/3\end{bmatrix} = \begin{bmatrix}1 \\ 1 \\ 2\end{bmatrix} - \begin{bmatrix}2/3 \\ -2/3 \\ 2/3\end{bmatrix} - \begin{bmatrix}5/6 \\ 5/3 \\ 5/6\end{bmatrix} = \begin{bmatrix}-1/2 \\ 0 \\ 1/2\end{bmatrix}$ \\
$C' = \frac{C}{||C||} = \frac{\begin{bmatrix}-1/2 \\ 0 \\ 1/2\end{bmatrix}}{\sqrt{1/4 + 0 + 1/4}} = \frac{\begin{bmatrix}-1/2 \\ 0 \\ 1/2\end{bmatrix}}{1/\sqrt{2}} = \begin{bmatrix}-1/\sqrt{2} \\ 0 \\ 1/\sqrt{2}\end{bmatrix}$

\noindent
The orthonormal set of vectors is:
\[
\left\{
\begin{bmatrix}
\frac{1}{\sqrt{3}} \\
-\frac{1}{\sqrt{3}} \\
\frac{1}{\sqrt{3}}
\end{bmatrix},
\begin{bmatrix}
\frac{1}{\sqrt{6}} \\
\frac{2}{\sqrt{6}} \\
\frac{1}{\sqrt{6}}
\end{bmatrix},
\begin{bmatrix}
-\frac{1}{\sqrt{2}} \\
0 \\
\frac{1}{\sqrt{2}}
\end{bmatrix}
\right\}
\]

%
% Question 6 %%%%%%%%%%%%%%%%%%%%%%%%%%%%%%%%%%%%%%%%%%%%%%%%%%%%%%%%%%%%%%%%%
\vfill\break
\section*{Question 6} (10 points) Consider a vector $a_e = \begin{bmatrix}3 & 1\end{bmatrix}$ defined on the natural basis.  Now, given the new basis defined by a matrix $B = \begin{bmatrix}2 & 2 \\ -2 & 1\end{bmatrix}$ how is the vector $a_e$ expressed in the new basis?

\noindent
To express the vector $a_e = \begin{bmatrix}3 & 1\end{bmatrix}$ in the new basis defined by the matrix $B = \begin{bmatrix}2 & 2 \\ -2 & 1\end{bmatrix}$, we need to find the coordinates of $a_e$ in the new basis. This can be done by solving the equation:
\[
a_e = B a_b
\]
where $a_b$ is the vector in the new basis. To find $a_b$, we can use the inverse of the matrix $B$:
\[
a_b = B^{-1} a_e
\]

\noindent
First, we need to find the inverse of the matrix $B$:
\[
B = \begin{bmatrix}2 & 2 \\ -2 & 1\end{bmatrix}
\]
The determinant of $B$ is:
\[
\det(B) = (2)(1) - (2)(-2) = 2 + 4 = 6
\]
The inverse of $B$ is:
\[
B^{-1} = \frac{1}{\det(B)} \begin{bmatrix}1 & -2 \\ 2 & 2\end{bmatrix} = \frac{1}{6} \begin{bmatrix}1 & -2 \\ 2 & 2\end{bmatrix} = \begin{bmatrix}\frac{1}{6} & -\frac{1}{3} \\ \frac{1}{3} & \frac{1}{3}\end{bmatrix}
\]
\noindent
Now, we can find the vector $a_b$ in the new basis:
\[
a_b = B^{-1} a_e = \begin{bmatrix}\frac{1}{6} & -\frac{1}{3} \\ \frac{1}{3} & \frac{1}{3}\end{bmatrix} \begin{bmatrix}3 \\ 1\end{bmatrix} = \begin{bmatrix}\frac{1}{6} \cdot 3 + (-\frac{1}{3}) \cdot 1 \\ \frac{1}{3} \cdot 3 + \frac{1}{3} \cdot 1\end{bmatrix} = \begin{bmatrix}\frac{1}{2} - \frac{1}{3} \\ 1 + \frac{1}{3}\end{bmatrix} = \begin{bmatrix}\frac{1}{6} \\ \frac{4}{3}\end{bmatrix}
\]
\noindent
The vector $a_e = \begin{bmatrix}3 & 1\end{bmatrix}$ expressed in the new basis defined by the matrix $B = \begin{bmatrix}2 & 2 \\ -2 & 1\end{bmatrix}$ is:
\[
a_b = \begin{bmatrix}\frac{1}{6} \\ \frac{4}{3}\end{bmatrix}
\]

%
% Question 7 %%%%%%%%%%%%%%%%%%%%%%%%%%%%%%%%%%%%%%%%%%%%%%%%%%%%%%%%%%%%%%%%%
\vfill\break
\section*{Question 7} (10 points) Find all the eigenvalues and eigenvectors of $A = \begin{bmatrix}1 & 2 \\ 4 & 3\end{bmatrix}$.

To find the eigenvalues and eigenvectors of the matrix \( A = \begin{bmatrix}1 & 2 \\ 4 & 3\end{bmatrix} \), we follow these steps:
% Step 1: Find the Eigenvalues
The eigenvalues \(\lambda\) of matrix \( A \) are found by solving the characteristic equation:
\[
\det(A - \lambda I) = 0
\]
where \( I \) is the identity matrix.
% Step 2: Compute the Characteristic Polynomial
\[
A - \lambda I = \begin{bmatrix}1 - \lambda & 2 \\ 4 & 3 - \lambda\end{bmatrix}
\]
The determinant of \( A - \lambda I \) is:
\[
\det(A - \lambda I) = (1 - \lambda)(3 - \lambda) - (2)(4) = (3 - \lambda - 3\lambda + \lambda^2) - 8 = \lambda^2 - 4\lambda - 5
\]
\[
\lambda^2 - 4\lambda - 5 = 0
\]
\[
\lambda = \frac{4 \pm \sqrt{16 + 20}}{2} = \frac{4 \pm 6}{2}
\]
The eigenvalues are:
\[
\lambda_1 = 5 \quad \text{and} \quad \lambda_2 = -1
\]
% Step 4: Find the Eigenvectors
For each eigenvalue, we find the corresponding eigenvector by solving \((A - \lambda I)v = 0\).

For \(\lambda_1 = 5\):
\[
A - 5I = \begin{bmatrix}1 - 5 & 2 \\ 4 & 3 - 5\end{bmatrix} = \begin{bmatrix}-4 & 2 \\ 4 & -2\end{bmatrix}
\]
The system of equations is:
\[ -4x + 2y = 0 \quad \]
\[ \quad 2x = y \]
The eigenvectors:
\[
v_1 = \begin{bmatrix}x \\ 2x\end{bmatrix} = x \begin{bmatrix}1 \\ 2\end{bmatrix}
\]
For \(x = 1\), we get:
\[
v_1 = \begin{bmatrix}1 \\ 2\end{bmatrix}
\]

For \(\lambda_2 = -1\):

\[
A - (-1)I = \begin{bmatrix}1 - (-1) & 2 \\ 4 & 3 - (-1)\end{bmatrix} = \begin{bmatrix}2 & 2 \\ 4 & 4\end{bmatrix}
\]
\[ 2x + 2y = 0 \quad \]
\[ \quad x = -y \]
Thus, the eigenvector can be written as:
\[
v_2 = \begin{bmatrix}x \\ -x\end{bmatrix} = x \begin{bmatrix}1 \\ -1\end{bmatrix}
\]
For \(x = 1\), we get:
\[
v_2 = \begin{bmatrix}1 \\ -1\end{bmatrix}
\]
The eigenvectors are:
\[
v_1 = \begin{bmatrix}1 \\ 2\end{bmatrix} \quad \text{and} \quad v_2 = \begin{bmatrix}1 \\ -1\end{bmatrix}
\]

The eigenvalues of the matrix \( A \) are \(\boxed{5}\) and \(\boxed{-1}\), with corresponding eigenvectors \(\boxed{\begin{bmatrix}1 \\ 2\end{bmatrix}}\) and \(\boxed{\begin{bmatrix}1 \\ -1\end{bmatrix}}\), respectively.

\vfill\break\noindent
\section*{Question 8} (5 Points) Find the best fit line using least squares approximation for the points (1, 2), (2, 1), (3, 1) in $R^2$

\textbf{Solution}
To find the best fit line using least squares approximation for the points \((1, 2)\), \((2, 1)\), and \((3, 1)\) in $R^2$, we follow these steps:

Compute the sums needed for the calculations:
\begin{enumerate}
      \item $\sum x_i = 1 + 2 + 3 = 6$
      \item $\sum y_i = 2 + 1 + 1 = 4$
      \item $\sum x_i y_i = (1)(2) + (2)(1) + (3)(1) = 2 + 2 + 3 = 7$
      \item $\sum x_i^2 = 1^2 + 2^2 + 3^2 = 1 + 4 + 9 = 14$
\end{enumerate}
The best fit line is given by the equation \( y = mx + b \), where the slope \( m \) and the y-intercept \( b \) are calculated as follows:

\noindent
The slope \( m \) is calculated using the formula:
\[
m = \frac{n \sum x_i y_i - \sum x_i \sum y_i}{n \sum x_i^2 - (\sum x_i)^2}
\]
\[
m = \frac{3 \times 7 - 6 \times 4}{3 \times 14 - 6^2} = \frac{21 - 24}{42 - 36} = \frac{-3}{6} = -0.5
\]

\noindent
The y-intercept \( b \) is calculated using the formula:
\[
b = \frac{\sum y_i - m \sum x_i}{n}
\]
Substituting the values:
\[
b = \frac{4 - (-0.5) \times 6}{3} = \frac{4 + 3}{3} = \frac{7}{3}
\]
Thus, the best fit line is:
\[
\boxed{y = -0.5x + \frac{7}{3}}
\]

\end{document}